\section{Introduction}
\label{sec:introduction}

Variational quantum algorithms and hybrid quantum-classical neural networks have emerged as candidate frameworks for extracting complex representations from high-dimensional data \cite{cerezo2021variational,schuld2021machine}. 
In medical imaging, multimodal volumetric acquisitions---such as multi-parametric magnetic resonance imaging (MRI) and computed tomography (CT)---provide complementary tissue properties indispensable for accurate diagnosis and tumour delineation \cite{menze2014multimodal,thummerer2023synthrad2023}. 
However, deploying parameterised quantum circuits (PQCs) within 3D medical vision pipelines exposes severe structural limitations. 
Contemporary noisy intermediate-scale quantum (NISQ) devices restrict processing to modest qubit counts ($N \le 12$) and shallow circuit depths ($L \le 5$) to avoid unmitigated decoherence and barren plateaus \cite{preskill2018quantum,mcclean2018barren,wang2021noise}.

When a high-dimensional 3D medical volume is compressed to fit the narrow input bandwidth of a quantum circuit, two competing failures typically manifest. 
First, if the network attempts to reconstruct voxel-level segmentations directly from low-dimensional quantum state measurements, fine boundary information is irreversibly destroyed---a phenomenon we define as the \textit{spatial reconstruction paradox}. 
Prior works frequently circumvented this issue by inserting high-capacity classical projection layers (often exceeding 200,000 parameters) after the quantum circuit. 
Such designs obscure whether reported performance gains originate from quantum representation utility or simply unconstrained classical capacity.

Second, when classical and quantum branches are fused through fixed, static combination ratios (e.g., 50/50 concatenation), optimization dynamics frequently degrade. 
Classical neural network layers propagate smooth, high-magnitude gradients, whereas quantum expectation measurements exhibit higher variance and susceptibility to hardware noise. 
In unregularized gating schemes, the routing mechanism rapidly degenerates to 100\% classical processing and 0\% quantum processing---a failure mode termed \textit{classical collapse}. 
Conversely, rigid 100\% quantum processing forces all representations through noisy, restricted quantum bottlenecks, degrading performance when input dimensionality or modality complexity increases.

In this work, we propose \textbf{Q-FUSION}, an adaptive classical-quantum fusion framework designed to operate under strict parameter and quantum-resource constraints. 
Rather than claiming unconditional quantum computational advantage, we investigate a concrete scientific hypothesis: 
\textit{Can a dynamic, trainable gating router learn to allocate computational representations between classical and quantum pathways to preserve complementary multimodal information, while spatial modulation shields volumetric reconstruction from severe quantum bottleneck constraints?}

The contributions of this study are structured as follows:
\begin{enumerate}
    \item \textbf{Resolution of Classical Collapse via Annealed Routing:} We introduce a piecewise linear annealed routing penalty that explicitly enforces pathway exploration during early training epochs before allowing the gating scalar $\alpha(x) \in [0, 1]$ to settle on an equilibrium allocation.
    \item \textbf{Quantum-Conditioned Spatial Modulation:} To bypass the spatial reconstruction paradox without high-capacity unflattening layers, we implement Feature-wise Linear Modulation (FiLM) \cite{perez2018film}. Compact 8-dimensional quantum state measurements modulate classical 3D feature maps through affine scaling and shifting parameters, preserving spatial geometry in the classical CNN while incorporating global quantum conditioning.
    \item \textbf{Strict Capacity-Matched Classical Controls:} We design parameter- and capacity-matched classical surrogate transformations possessing identical weight counts (24 parameters) and identical value bounds ($[-1, 1]$), ensuring exact 0.0000\% capacity variance between classical-quantum and classical-classical architectures.
    \item \textbf{Empirical Multi-Modality Scaling on BraTS2021:} On a standardized 150-patient cohort from BraTS2021 \cite{baid2021rsna}, we evaluate scaling across $M \in \{1, 2, 3, 4\}$ physical MRI sequences (T1, T1ce, T2, FLAIR). We demonstrate that at $M=4$, classical pooling suffers a 4.41\% DSC degradation due to feature dilution, whereas adaptive routing maintains representation stability (+4.79\% DSC over classical baselines) through automated routing rebalancing.
    \item \textbf{NISQ Hardware Noise Characterization:} We benchmark open-system gate depolarizing noise ($p \in [0.000, 0.050]$) on density-matrix simulators, demonstrating that adaptive routing provides structural buffering against physical hardware decoherence.
\end{enumerate}
